\section{Discussion and limitations}
\label{sec:discussion}

\paragraph{What is new and what is not} That SR models trained on bicubic downsampling degrade on compressed inputs is known for natural images. What this study adds is a measurement in a regime and a domain where it had not been made: inputs of a few dozen pixels of a biometric trait, where the compression level at which published efficient models stop being useful (about quality 85) falls between the two levels that real images of the domain actually have; the observation that a generic real-world pipeline is not a remedy there; and the observation that, once the degradation is repaired, a \ParamsSpan{}M-parameter network suffices to exceed every published model we tested in fidelity, including one \ParamsRatio{} times larger.

\paragraph{Limits of the claim} (i) The model is better than the published fidelity-oriented models and than bicubic interpolation on compressed inputs, in fidelity and in LPIPS. It is not better on clean inputs, and large GAN-trained models keep a lower LPIPS at a much lower fidelity and a much higher cost. (ii) Most of the gain over a fixed-quality JPEG degradation is explained by the distribution of compression levels (Section~\ref{sec:ablation}); we do not claim that estimating noise and blur is necessary. (iii) Fidelity results rely on synthesised inputs. The in-the-wild datasets have ground truth only at the smallest sizes, and the evaluation on real small images is through identification and viewer preference, not fidelity. (iv) AMI was photographed indoors and has 100 subjects; each fold is trained with one seed, so the spread across folds stands in for the spread across seeds. (v) The realism classifier is evaluated on \RealTestSubj{} held-out subjects; an accuracy near 50\% shows that this classifier cannot separate the two sets, not that the distributions coincide. (vi) The blur of real small images is matched only through a sharpness statistic.

\paragraph{Practical guidance} For small ear images that have been through JPEG compression, a published efficient SR model should not be assumed to improve on interpolation. Fine-tuning on ear images is not sufficient if the degradation is left unchanged, and a generic real-world degradation is not a substitute for one that covers the compression levels of the target images. When the fine-tuning data come from a controlled acquisition, photometric augmentation and a brightness stress test should be part of the procedure.

\section{Conclusion}
On small ear images, published efficient SR models fall below bicubic interpolation in fidelity under the compression that real images of the domain carry, and no choice among sixteen architectures changes this. Repairing the training degradation does: a real-time network fine-tuned with a degradation measured from real small ear images, under a recipe that keeps fine-tuning stable outside the laboratory data, exceeds bicubic interpolation and the published models on three datasets at no additional cost. The benchmark, the measured degradation and the baselines are released so that architectures designed for such inputs can be compared against a model that is trained correctly.
